% Generated draft: requires mathematical and source review.
\documentclass[11pt]{article}
\usepackage[margin=25mm]{geometry}
\usepackage{fontspec}
\setmainfont{Latin Modern Roman}
\usepackage{amsmath,amssymb,mathtools,unicode-math}
\setmathfont{Latin Modern Math}
\usepackage{enumitem,needspace,xurl,hyperref}
\setlength{\parindent}{0pt}
\setlength{\parskip}{5pt}
\setlength{\emergencystretch}{3em}
\newcommand{\N}{\mathbb{N}}
\newcommand{\Z}{\mathbb{Z}}
\newcommand{\Q}{\mathbb{Q}}
\newcommand{\R}{\mathbb{R}}
\newcommand{\C}{\mathbb{C}}
\newcommand{\F}{\mathbb{F}}
\newcommand{\A}{\mathbb{A}}
\newcommand{\PP}{\mathbb{P}}
\newcommand{\E}{\mathbb{E}}
\newcommand{\eps}{\varepsilon}
\newcommand{\dd}{\,\mathrm{d}}
\begin{document}
\section*{UnsolvedMath ID: 11000081}
\textbf{Source problem code:} catalogue entry
\subsection{Definitions and mathematical statement}
Fix an integer $g\ge2$ and let $\operatorname{Mod}_g$ be the mapping class group of a closed orientable genus-$g$ surface. For a pseudo-Anosov mapping class $f$, let $\lambda(f)>1$ denote its stretch factor. With the usual Teichm\"uller metric normalization, its translation length and associated closed-geodesic length are $\log\lambda(f)$.

The preceding local source item, ID 11000080, defines $(k_1,\ldots,k_r)$ as prong counts at the $r$ singularities of a stable foliation, not orders of zeros of a differential. For a closed unmarked surface, $k_i\ge3$ and the Poincar\'e--Hopf constraint is
\[
 \sum_{i=1}^r(k_i-2)=4g-4.
\]
Call a tuple admissible when it is realizable in the intended foliation setting; this is not an assertion that the displayed numerical condition alone is sufficient. The source attributes pseudo-Anosov realization of every admissible tuple to Masur--Smillie [MS]. No orientability or connected-stratum-component condition is imposed in the question.

Define the minimum-length infimum for those singularity data by
\[
 \lambda_g(k_1,\ldots,k_r)=
 \inf\{\log\lambda(f):f\in\operatorname{Mod}_g\text{ pseudo-Anosov, with stable foliation having prong counts }(k_1,\ldots,k_r)\}.
\]
Although this symbol uses $\lambda$, it denotes a logarithmic length, not the stretch factor itself.

\textbf{Problem 7.5.} For every fixed $g\ge2$ and every admissible tuple, give upper and lower bounds for $\lambda_g(k_1,\ldots,k_r)$ in terms of the genus and singularity data. The source interprets this as bounding the shortest Teichm\"uller loop in the corresponding substratum in moduli space, obtained from the relevant cotangent-bundle substratum. The defining condition is on the invariant foliation; merely specifying that the projected loop lies in a projection of a stratum can forget this data, so the infimum formula remains the precise target.

Writing
\[
 L(\operatorname{Mod}_g)=\inf_{f\text{ pseudo-Anosov}}\log\lambda(f),
\]
the source notes
\[
 L(\operatorname{Mod}_g)=\min_{\text{admissible }(k_1,\ldots,k_r)}
 \lambda_g(k_1,\ldots,k_r),
\]
where only finitely many singularity tuples occur for fixed $g$. The problem concerns bounds for every specified tuple rather than only this overall minimum.
\subsection{Short English statement}
For each admissible stable-foliation prong pattern in genus at least two, bound the infimum of pseudo-Anosov logarithmic stretch factors with exactly that pattern. Taking the minimum over patterns recovers the unrestricted shortest Teichm\"uller-loop length.
\subsection{Sources}
\begin{itemize}
\item Local UnsolvedMath registry, permanent ID 11000081, source code catalogue entry; no individual proposer or proposal year supplied.
\item Farb, edited volume \emph{Problems on Mapping Class Groups and Related Topics}, Problem 7.5, PDF page 58; \url{https://www.math.uchicago.edu/~farb/papers/mcgbook.pdf}. The registry records PDF-font-structure extraction and access on 2026-07-29; the inventory distinguishes 329 labeled headings by source page.
\item Adjacent local ID 11000080 supplies the prong-count convention, finite-tuple observation, and Masur--Smillie realization attribution [MS]. No full bibliography entry for [MS] is supplied here.
\end{itemize}
\end{document}
